\section{Related Work}

Research on long-context KV management has taken several distinct directions. Quantization methods preserve the full cache structure while reducing the number of stored bits per vector. KIVI uses asymmetric low-bit treatment for keys and values \citep{liu2024kivi}, and later work has examined other quantized or shared-cache variants such as CommonKV \citep{wang2025commonkv}. These approaches are attractive because they retain all positions, but they do not by themselves reduce the number of attended entries.

Eviction and retrieval methods attack the compute side more directly. H2O retains heavy hitters \citep{zhang2023h2o}; Ada-KV reallocates budgets adaptively across heads \citep{feng2024adakv}; SnapKV uses an observation window to select positions that appear useful for future queries \citep{li2024snapkv}; Quest scores pages or blocks with a query-aware rule \citep{tang2024quest}; and A2ATS argues that value information also matters for token selection \citep{guo2024a2ats}. These methods can reduce attention cost, but they are primarily heuristic or predictive. Their selection rules are not local error certificates in the sense studied here.

Residual and reference-based compression lines are closer in spirit to RACK-KV because they exploit temporal structure in historical states. DeltaKV compresses through long-range residual similarity \citep{hao2026deltakv}, OjaKV uses online low-rank structure \citep{zhu2025ojakv}, EchoKV reconstructs from similarity-based summaries \citep{ji2026echokv}, and probabilistic language tries provide a sequential predictive viewpoint on KV compression \citep{magarshak2026pltries}. These methods show that historical KV states contain exploitable structure, but they do not simultaneously provide an independently decodable random-access format together with a formal skip certificate.

System work on memory organization provides a different piece of the puzzle. PagedAttention addresses efficient paging and allocation for serving \citep{kwon2023pagedattention}. InfLLM and related context-memory systems restructure access to long contexts \citep{xiao2024infllm}. DeepSeek-V2 reduces KV demand partly through architectural changes rather than cache compression alone \citep{deepseekai2024deepseekv2}. These systems motivate the importance of access patterns and bandwidth, but they do not offer the same local theorem for attention-output error under skipping.

\begin{table}[t]
  \centering
  \small
  \begin{threeparttable}
\begin{tabularx}{\linewidth}{@{}lXcccc@{}}
\toprule
Method & Core idea & Removes history & Query-aware & Bound & Status here \\
\midrule
Full KV & Exact cache & No & No & Exact & Reference \\
Uniform INT8 & Uniform KV quantization & No & No & None & Baseline \\
KIVI & Asymmetric low-bit quantization & No & No & None & KIVI-style \\
SnapKV & Window-based token selection & Yes & Yes & None & SnapKV-style \\
Quest & Query-aware page retrieval & Yes & Yes & None & Quest-style \\
DeltaKV & Residual-style compression & No & No & None & Literature \\
RACK-KV & Block-local compression plus certificate & Optional & Yes & Local & Method \\
\bottomrule
\end{tabularx}
\end{threeparttable}

  \caption{Compact comparison with representative long-context methods. The project baselines labeled \enquote{-style} are source-inspired approximations rather than exact reproductions of the original systems.}
  \label{tab:related}
\end{table}

\paragraph{Provisional corpus-bounded novelty assessment.}
No fully verified work in the frozen corpus was found that combines sequence-aware or reference-based KV compression, independently decodable random-access compressed blocks, and a formal safe-skipping certificate in one inference-time system. This statement is deliberately narrow. It is not a claim of absolute worldwide novelty, it is limited to the frozen July 2026 corpus, and individual ingredients of RACK-KV overlap prior lines of work. DeltaKV is the strongest direct overlap on residual-style compression; probabilistic language tries are the main threat on sequential predictive coding; and Quest is adjacent on query-aware block scoring, although without the same certified local output bound.
